% ============================================================================
% ADDENDUM TO LIMITATIONS SECTION
% ============================================================================
% Add this to the Limitations section (Section 6 or Appendix)

\paragraph{Regime-Dependent Hyperparameter Effects.}
Our sensitivity analysis (Section~\ref{sec:sensitivity_analysis}) uses seed 42 as the primary illustrative example, which represents a warmup-dominant regime where Corralling maintains the semantic transfer expert throughout evaluation. Multi-seed validation (Appendix~\ref{app:multiseed_sensitivity}) reveals this regime occurs in approximately 33\% of data orderings, with the remaining 67\% exhibiting tabula rasa-dominant behavior where Corralling abandons semantic transfer in favor of cold-start exploration.

This regime-dependence is a \textit{feature, not a limitation}. It demonstrates Corralling's adaptive capacity: when warmup priors match the data distribution (warmup-dominant regime), hyperparameter choice ($\neff$) significantly impacts performance (+4.6\% for optimal vs strong priors). When priors mismatch the data (tabula rasa-dominant regime, characterized by 71.5\% ties and low task variance), Corralling correctly detects the over-confidence and switches strategies, rendering hyperparameter choice irrelevant (0\% effect).

The regime switching mechanism depends on early-phase performance (t=0-300) of warmup priors relative to cold-start exploration. In production deployment, the specific regime observed will depend on:
\begin{enumerate}[leftmargin=*, itemsep=2pt]
    \item \textbf{Prior Quality}: Whether warmup priors reflect the actual test distribution
    \item \textbf{Task Heterogeneity}: Proportion of prompts where model preferences are decisive vs tied
    \item \textbf{Data Ordering}: Which task types appear in early evaluation (though Corralling adapts within 100-200 steps)
\end{enumerate}

For practitioners, this implies:
\begin{itemize}[leftmargin=*, itemsep=2pt]
    \item \textbf{Monitoring}: Track expert selection frequencies ($\sim$30-40\% warmup expected based on our data). Persistent 0\% or 100\% may indicate distribution shift or prior quality issues.
    \item \textbf{Hyperparameter Tuning}: Limited value for $\neff$ optimization (expected impact $\sim$1.5\% overall). Focus on prior quality and learning rate ($\eta$) selection instead.
    \item \textbf{Robustness Validation}: Regime-dependent effects require stratified analysis. When reporting aggregate performance, specify whether results reflect warmup-dominant, tabula rasa-dominant, or mixed regimes.
\end{itemize}

\paragraph{Generalizability of Regime Frequencies.}
The 33\%/67\% warmup/tabula rasa split observed in our experiments reflects the specific characteristics of our test distribution (LMSYS Chatbot Arena battles with GPT-5.1 vs GPT-4-Turbo vs Mixtral-8x7B). Different model portfolios or task distributions may exhibit different regime frequencies. For instance, if warmup priors are high-quality (trained on well-matched data) and the task distribution exhibits high heterogeneity, warmup-dominant regimes may occur more frequently (>50\%). Conversely, if priors are misspecified or task variance is low (many ties), tabula rasa regimes may dominate (>80\%).

The adaptive behavior generalizes: Corralling will always monitor expert performance and shift probability mass toward better-performing strategies. What varies across deployments is the \textit{frequency} of each regime, not the \textit{mechanism} of regime switching.

% ============================================================================
% END ADDENDUM
% ============================================================================
